\chapter{Обзор литературы}\label{ch:ch1}

В этой главе приведен обзор литературы и методов, используемых в данной работе. Первый подраздел посвящен обзору метода молекулярной динамики от основ до конкретных приожений, затронутых  в данной работе. Также приведены сведения и разобраны примеры применения потенциалов межатомного взаимодействия, применявшихся в этой работе. Во втором подразделе описаны фазовые переходы первого рода, история их исследования и современные представлениям о них. Третий раздел посвящен поверхностным явлениям, фазовому равновесию и методам их изучения методом молекулярной динамики. В последнем разделе этой главы описаны процессы переноса и коэффициент термической аккомодации.

Настоящая глава формирует теоретическую и методологическую базу диссертационного исследования. Её цель — систематизировать ключевые концепции и инструменты, необходимые для моделирования и анализа процессов фазовых превращений в рамках молекулярно-динамического подхода. 

В первом разделе даётся детальное описание метода молекулярной динамики (МД). Особое внимание уделяется физическим основам и математическому аппарату МД, а также обзору потенциалов межатомного взаимодействия, используемых в работе, поскольку их адекватный выбор является залогом достоверности получаемых результатов.

Второй раздел посвящён классической и современной проблематике фазовых переходов первого рода. Рассматривается их термодинамическая классификация, история изучения и фундаментальный механизм зарождения новой фазы — нуклеация.

Третий раздел фокусируется на поверхностных явлениях и фазовом равновесии. Подробно разбираются методы вычисления поверхностного натяжения — ключевого параметра, определяющего энергетику нуклеации, непосредственно из данных МД-моделирования.

Наконец, в четвёртом разделе анализируются процессы переноса. Изучается роль диффузии в кинетике роста зародышей и вводится коэффициент термической аккомодации как фундаментальная характеристика, связывающая атомистические процессы на межфазной границе с макроскопическим описанием тепло- и массообмена. 

\section{Метод молекулярной динамики}

\subsection{Общие сведения}

Метод молекулярной динамики заключается в решении уравнений Ньютона для всех $N$ частиц системы:

\begin{align}
	m_i\frac{\overrightarrow{d^2r_i}}{dt^2} = \overrightarrow{F_i}(r_1, r_2, ... r_n),
\end{align}

где $m_i$ \--- масса частицы с номером $i$, $r_i$ \--- её координата, а $F_i$ \--- скорость. Сила рассчитывается как  градиент потенциальной энергии, взятый со знаком минус:

\begin{align}
	\overrightarrow{F_i}(r_1, r_2, ... r_n) = -\overrightarrow{\nabla}U_i(r_1, r_2, ... r_n),
\end{align}

где $U_i$ \--- потенциальная энергия $i$-ого атома. Она определяется через ряд функций, называемых потенциалом взаимодействия. Потенциал межатомного взаимодействия определяет свойства системы. 

Далее для расчета положений атомов и их скоростей в следующий шаг по времени применяется схема Верле:

\begin{subequations}
    \begin{gather}
		\overrightarrow{r_i}(t+\Delta t) = \overrightarrow{r_i}(t) + \overrightarrow{v_i}(t) + \frac{\overrightarrow{F_i}(t)\Delta t^2}{2m_i},\\
		\overrightarrow{v_i}(t+\Delta t) = \overrightarrow{v_i}(t) + \frac{\overrightarrow{F_i}(t)+\overrightarrow{F_i}(t+\Delta t)}{2m_i}\Delta t.
    \end{gather}
\end{subequations}

Из-за наличия численных ошибок, энергия в молекулярно-динамической системе не сохраняется. Поскольку в природе частицы двигались бы по оптимальному пути с минимальной энергией, отклонения от него увеличивают энергию системы. Такие прибавки к энергии зависят от шага по времени. Однако, слишком малый шаг по времени требует большого времени расчетов. Следовательно, требуется соблюдать баланс и учитывать оба фактора.

В случае, если система состоит не только из атомов, но и молекул, схема Верле становится неэффективной. Необходимо, чтобы шаг по времени был на порядок меньше самых быстрых колебаний в системе. В случае наличия жестких связей или углов это время может быть на порядок меньше шага по времени, достаточного для межмолекулярных взаимодействий.

В таких случаях используется иерархическая схема rRESPA(reversible REference System Propagator Algorithm)\cite{tuckerman_reversible_1992}. Это иерархическая схема, позволяющая рассчитывать координаты и скорости учитывая каждое взаимодействие с разной частотой. Например, можно сделать шаг по времени для связей в 10 раз короче, чем для межмолекулярных взаимодействий. Это достигается разложением оператора Лиувилля на сумму двух слагаемых, отвечающих за разные потенциалы:

\begin{align}
	iL = iL_1 + iL_2,
\end{align}

Для такого разложения справедливо:

\begin{align}
	e^{i(L_1+L_2)t} = (e^{i(L_1 + L_2)t/N_t})^{N_t} = (e^{i L_1 \Delta t /2} e^{i L_2 \Delta t} e^{i L_1 \Delta t /2})^{N_t} + O(t^3/P^2),
\end{align}

где $\Delta t = t/N_t$ \--- шаг по времени, а $t$ и $N_t$ \--- время моделирования и число шагов соответственно. Тогда оператор эволюции примет вид:

\begin{align}
	 e^{i L \Delta t} = e^{i L_1 \Delta t /2} e^{i L_2 \Delta t} e^{i L_1 \Delta t /2}.
\end{align}
  
Такая схема позволяет интегрировать "быстрый" процесс с шагом в $N$ раз меньшим, чем "медленный", что позволит достичь лучшей устойчивости модели и сохранения энергии.

Для начала моделирования требуется задать начальные условия. Два наиболее распротраненных метода задания начальных координат \--- случайное расположение атомов в заданной области и расположение их в узлах кристаллической решетки. У каждого из методов есть свои особенности. В первом случае атомы могут оказаться достаточно близко друг к другу, чтобы система оказалась нестабильной. Потому проводится процесс минимизации энергии системы до непосредственного моделирования. В случае расположения в узлах кристаллической решетки этого удается избежать сразу и даже решить проблему с циклическими молекулами, которые могут пересечься и не разъединиться в результате минимизации энергии. Однако, из-за этого система приобретает упорядоченность, которая может быть ей не свойственна. В таком случае потребуется долго моделировать систему, чтобы она вышла на равновесие.

Ещё одной важной деталью молекулярного моделирования являются граничные условия. Как правило, используются периодические. В этом случае создаются образы вычислительной ячейки периодически в каждом направлении. Таким образом, атомы, находящиеся на границе ячейки всё равно находятся в окружении других атомов как если бы моделировалась непрерывная среда. Такой вариант подходит для большинства систем. В некоторых случаях, как, например, в главе~\cref{ch:ch6} описана система, состоящая из уединенного кластера атомов железа и налетающего атома. Там используются фиксированные граничные условия, поскольку в случае периодических могло бы быть, например, несколько соударений одного и того же атома с кластером, что затруднило бы обработку данных.

\subsection{ Потенциалы межатомного взаимодействия}

Наиболее важной характеристикой, определяющей систему, является потенциал межатомного взаимодействия. Он определяет силы, действующие между атомами в вычислительной ячейке.

Самым первым использованным потенциалом является потенциал Леннард-Джонса:

\begin{align}
	U(\overrightarrow{r}) = 4 \epsilon (\frac{\sigma}{r}^{12}-\frac{\sigma}{r}^6),
\end{align}

где $\epsilon$ соответствует минимальной потенциальной энергии взаимодействия двух частиц, так называемой "глубине потенциальной ямы", а $\sigma$ \--- расстоянию, где силы отталкивания равны силам притяжения между двумя частицами, то есть потенциальная энергия проходит через 0.

Этот потенциал хорошо описывает свойства инертных газов, соответствует энергии взаимодействия двух наведенных диполей и очень прост в расчётах: для вычисления члена, отвечающего за отталкивание, нужно $(\sigma / r)^6$ возвести в квадрат. Также этот потенциал используется для расчета межмолекулярных взаимодействий в углеводородах. Например, в моделях TraPPE-UA, OPLS и COMPASS. 

Однако, не для всех систем подходит такое описание, и разрабатываются иные потенциалы межатомного взаимодействия. Например, в потенциале Букингема~\cite{buckingham_1938} изменен член, отвечающий за отталкивание, на экспоненциальный:

\begin{align}
	U(\overrightarrow{r}) = A e^{-B r} - \frac{C}{r}^6),
\end{align}

где $A, B, C$ \--- константы. 

Наряду с потенциалом Букингема в таких задачах, как взаимодействие атомов металлов с инертным газом, используется потенциал Морзе~\cite{}. В нем оба слагаемых заменены на экспоненциальные функции:

\begin{align}
	U(\overrightarrow{r}) = D [e^{-2 \alpha (r - r_0)} - 2 e^{-\alpha (r - r_0)}],
\end{align}

где $D$ \--- глубина потенциальной ямы, $r_0$ \--- расстояние между атомами, на котором находится минимум потенциальной энергии, $\alpha$ \--- коэффициент.

Все эти потенциалы межатомного взаимодействия учитывают только расстояния между парами атомов. Такой способ расчета потенциальной энергии плохо подходит, например, для металлов, для которых необходимо учитывать общую электронную плотность. Для этого существуют многочастичные потенциалы, наподобие EAM(embedded atom method), в которых взаимодействие между атомами зависит также от величины электронной плотности, в которую вносят вклад все соседние атомы. Эта величина, как и все остальные, необходимые для данного потенциала, задаются таблично. Таким образом, формула для EAM принимает вид:

\begin{align}
	U(\overrightarrow{r}) = U_{pair}(r) + U_{many}(\rho),
\end{align}

где $U_pair$ \--- энергия парного взаимодействия, $U_{many}$ \--- энергия, зависящая от электронной плотности.

Данный тип потенциалов получил широкое развитие, поскольку достаточно хорошо описывают поведение металлов. В частности, можно отметитпотенциал Финниса-Синклера, который хорошо описывает плавление металлов и отличается тем, что учитывает тип атомов, которыь е вносят вклад в электронную плотность:

\begin{align}
	U(\overrightarrow{r}) = U_{pair}(r) + U_{\alpha \beta}(\rho),
\end{align}

где $U_{\alpha \beta}$ \--- энергия многочастичного взаимодействия, зависящая от типов атомов.

Однако, взаимодействия между атомами могут быть не только за счет сил Ван-дер-Ваальса. Внутримолекулярные взаимодействия описываются сходным образом. Существует множество моделей, описывающих молекулы. Чаще всего в них энергия взаимодействия  атомов разбивается на энергии связей, углов и двугранных углов:

\begin{align}
	U(r) = U_{bond} + U_{angle} + U_{dihedral}
\end{align}

Существует множество способов описания этих взаимодействий. Например, для связей часто применяется гармонический потенциал:

\begin{align}
	U(r) = k (r - r_0)^2,
\end{align}

где $r$ \--- длина связи, $r_0$ \--- равновесная длина связи, $k$ \--- коэффициент жесткости. Следует отметить, что коэффициент жесткости, например, в углеводородах достаточно велик, что вынуждает рассчитывать связи по схеме RESPA с меньшим шагом по времени.

Аналогично рассчитываются энергии углов. В случае с двугранными углами используются, например, суммы Фурье:

\begin{align}
	U(\phi) = \sum_{i=1, m}[K_i(1+cos(n_i\phi - d_i))],
\end{align}

В случае модели COMPASS имеется вклад энергии связей в энергию углов и двугранных углов, что будет более подробно показано в главе~\cref{ch:ch2}.

В некоторых случаях, например, в модели воды TIP4P, применяются жесткие связи и углы. Это означает, что длины связей и величины углов остаются неизменными. Это достигается, например, с помощью алгоритма SHAKE. 

\subsection{Дальнодействие и производительность}

Очевидно, что выбор потенциала межатомного взаимодействия влияет на скорость проведения расчетов. Одним из основных факторов, влияющих на производительность, является радиус обрезания. Это расстояние, начиная с которого энергия взаимодействия между атомами приравнивается к нулю. Например, для потенциала Леннард-Джонса:

\[
U(r) = \left\{
\begin{alignedat}{2}
    &= 4 \epsilon (\frac{\sigma}{r}^{12}-\frac{\sigma}{r}^6), \quad &\text{eсли } r\geqslant r_{cut} \\
    &0, \quad & \text{eсли } r_{cut}<0
\end{alignedat}
\right.
\]

Таким образом удается достичь сложности алгоритма $O(Nlog(N))$, вместо $N^2$, где $N$ \--- число частиц. При этом возникает ещё и такой объект как список соседей. В него включаются все атомы, находящиеся ближе, чем некоторый радиус, необязательно равный радиусу обрезания, к рассматриваемому атому. Зачастую радиус для списка соседей выбирается больше радиуса обрезания, благодаря чему возникает скин-слой, атомы в котором находятся в списке соседей, но не взаимодействуют с рассматриваемым атомом. Это позволяет реже обновлять список соседей и исключить возможность появления в пределах радиуса обрезания атомов, взаимодействие с которыми не учитывается, что позволяет улучшить производительность.  

Очевидно, что такой метод провоцирует недооценку потенциальной энергии. Существует несколько методов её учесть. Следует заметить, что интегралы вышеописанных потенциалов на бесконечности сходятся:

\begin{align}
	\label{eq:integral}
	U_{tail} = \int_{r_{cut}}^{\inf} \rho U(r) 4 \pi r^2 dr,
\end{align}

где $\rho$ \--- плотность, $U(r)$ \--- рассматриваемый потенциал межатомного взаимодействия.

При помощи подобного интеграла рассчитываются дальнодействующие поправки к потенциалу межатомного взаимодействия. Это позволяет уменьшить ошибку, возникающую из-за наличия радиуса обрезания. 

У этого метода есть ряд недостатков. Например, он не подходит для гетерогенной среды. Если моделируется две или более фазы, плотности в них будут разными, что противоречит предположению, сделанному в формуле.

Для взаимодействий по закону Кулона такой метод также не является подходящим, поскольку интеграл выше~\cref{eq:integral} на бесконечности сходится условно. Потому используются, например, потенциалы для плазмы, такие как потенциал Юкавы:

\begin{align}
    U(r) = - \frac{Z_1 Z_2 q_e^2}{k_1} e^{-k_2r},
\end{align}

где $k_1, k_2$ \--- некоторые константы, $Z_1, Z_2$ \--- зарядовые числа взаимодействующих частиц, $q_e$ \--- заряд электрона. Этот потенциал описывает экранирование, которое возникает в плазме, за счет обратного экспоненциального множителя. Такой потенциал является интегрируемым на бесконечности, потому для него возможно рассчитать поправки, компенсирующие отсутствие дальнодействующей части.

Ещё одним способом учесть дальнодействие в системе является метод Эвальда, который иногда, например, в GROMACS, реализован и для потенциалов, обычно рассчитываемых через поправки. Этот метод заключается в помещении заряда в оболочку противоположного знака. Далее происходит разложение в ряд Фурье для заряда в оболочке, которая задается формулой:

\begin{align}
    \rho_{Gauss}(r) = - q_i (\alpha / \pi)^{3/2} e^{-\alpha r^2},
\end{align}


где $\rho_{Gauss}$ \--- электронная плотность, задаваемая гауссианой, $q_i$ \--- экранируемый заряд, $\alpha$ \--- коэффициент, определяемый из соображений вычислительной эффективности.

В таком случае дальнодействующая часть будет равна:

\begin{align}
    U_1 = \frac{1}{2}\sum_{i}q_i \phi(r_i) = \frac{1}{2V}\sum_{k \neq 0} \frac{4\pi}{k^2}\rho(\overrightarrow{r})e^{-k^2/4\alpha},
\end{align}

где $\phi(r_i)$ \--- значение потенциала кулоновского потенциала в точке, где располагается $i$-ый атом, $k = 2\pi/L$, где $L$ длина кубической вычислительной ячейки.

Однако, это выражение включает взаимодействие точечного заряда и окружающей его гауссианы. Потому требуется вычесть следующее выражение:

\begin{align}
    U_{self} = (\alpha/\pi)^{1/2} \sum_{i = 1}^N q_i^2.
\end{align}

Таким образом, выражение для полной энергии кулоновского взамиодействия, включающее в себя и короткодействующую часть, примет вид:

\begin{equation}
    U_{self} = \frac{1}{2}\sum_{i}q_i \phi(r_i) = \frac{1}{2V}\sum_{k \neq 0} \frac{4\pi}{k^2}\rho(\overrightarrow{r})e^{-k^2/4\alpha} - (\alpha/\pi)^{1/2} \sum_{i = 1}^N q_i^2 + \frac{1}{2}\sum_{i \neq j}^N \frac{q_i q_j erfc(\alpha r_{ij})}{r_{ij}}.
\end{equation}

Следует заметить, что это требующая значительных ресурсов по сравнению с обычным расчетом дальнодействующих поправок процедура, потому был найден способ её ускорить. Для этого создается решетка, на которую проецируются имеющиеся в системе заряды. Этот метод называется particle-particle particle-mesh(PPPM).

Ещё один метод учета дальнодействия в системе показан в работе (Смирнов). Этот метод основан на моделировании системы с несколькими радиусами обрезания и экстраполяции полученных результатов на бесконечность. Такой метод показал хорошее согласие с экспериментом и достаточно высокую по сравнению с расчетом методом Эвальда производительность.

\subsection{Термостаты и баростаты}

Моделирование методом молекулярной динамики позволяет создавать системы, фиксируя различные условия. Самым простым вариантом является моделирование закрытого сосуда в теплоизолирующей оболочке, NVE-ансамбль. В этом случае фиксированны количество частиц, объем и энергия. 

Существует возможность зафиксировать температуру при моделировании, получив таким образом NVT-ансамбль, в котором происходит обмен энергией с окружающей средой, но не массой. Для этого потребуется усложнить алгоритм, добавив в систему термостат. Существует три самых распространенных вида термостатов: Берендсена, Ланжевена и Нозе-Гувера.

Термостат Берендсена является простейшим из перечисленных:

\begin{align}
    \frac{dT}{dt} = \frac{T-T_0}{\tau},
\end{align}

где $T$ и $T_0$ \--- температура системы и заданная температура соответсвенно, $t$ \--- время, $\tau$ \--- параметр, определяющий силу термостата. Его влияние заключается в следующем: каждая молекула системы, в зависимости от её текущей скорости, подвергается корректировке с использованием заданного коэффициента термализации. Это позволяет "отбирать" или "добавлять" энергию в систему, чтобы поддерживать температуру на заранее определённом уровне. 

Одним из недостатков этого термостата является нарушение динамики за счет изменения скоростей не вполне естественным образом. Также в результате работы этого термостата не устанавливается корректное распределение Гиббса.

Термостат Ланжевена также иногда называется стохастическим за счет добавления в него случайной силы:

\begin{align}
    m\frac{dv}{dt} = \frac{dU}{dr} - \lambda v + \xi(t),
\end{align}

где $v$ \--- скорость, $\lambda$ \--- коэффициент трения, $\xi(t)$ \--- некоторая функция для создания шума в системе. Коэффициент трения приводит к экспоненциальному уменьшению температуры, а шум подбирается так, чтобы приближаться к нужной температуре. Стохастичность уравнений движения, однако, принуждает переходить в систему отсчета, связанную с центром масс системы, чтобы избежать случая, когда частицы перестают двигаться относительно друг друга, но приобретают значительные сонаправленные скорости. Также этот термостат влияет на значения коэффициентов переноса.

Одним из наиболее широко используемых термостатов является термостат Нозе-Гувера, основанный на внесении поправок в лагранжиан системы. Новый лагранжиан принимает вид:

\begin{align}
    L_{Nose} = \sum_{i=1}^{N} \frac{1}{2} m_i s^2 \dot{\vec{r_i^2}} - U(\vec{r}^N) _ \frac{1}{2}Q \dot{s}^2 - \frac{3N}{kT} ln s,
\end{align}

где $s$ \--- параметр, характеризующий силу термостата, $Q$ \--- масса, которая к нему относится. 

В таком случае выражения для импульсов меняют вид:

\begin{align}
    \vec{p_i} = \frac{\partial L}{\partial \dot{\vec{r_i}}} = m_i s^2 \dot{\vec{r_i}}
\end{align}

Для новой координаты также определяется импульс:

\begin{align}
    p_s = \frac{\partial L}{\partial \dot{s}} = Q \dot{s}
\end{align}

Тогда гамильтониан системы примет вид:

\begin{equation}
    H_{Nose} = \sum_{i=1}^N \frac{\vec{p_i^2}}{2m_is^2} + U(\vec{r}^N) + \frac{p_s^2}{2Q} + \frac{3N ln s}{kT}
\end{equation}

Можно представить новые уравнения движения как уравнения движения с переменным шагом по времени. $s$ служит коэффициентом масштабирования между реальным шагом и тем, с которым происходит интегрирование. Сами уравнения движения примут вид:

\begin{subequations}
	\begin{gather}
		\frac{d\vec{r_i}}{dt} = \frac{\partial H_{Nose}}{\partial \vec{p_i}} = \frac{\vec{p_i}}{m_i s^2} \\
		\frac{d\vec{p_i}}{dt} = -\frac{\partial H_{Nose}}{\partial \vec{r_i}} = \frac{\partial U(\vec{r}^N)}{\partial \vec{r_i}} \\
		\frac{ds}{dt} = \frac{\partial H_{Nose}}{\partial p_s} = \frac{p_s}{Q} \\
		\frac{d p_s}{dt} = -\frac{\partial H_{Nose}}{\partial s} = \frac{1}{s} \sum_{i=1}^N (\frac{\vec{p_i}^2}{m_i s^2} - \frac{3N}{kT}) 
	\end{gather}
\end{subequations}

Недостатком термостата Нозе-Гувера являются нефизичные флуктуации температуры. Для уменьшения их влияния используется цепочка из таких термостатов, применяемых последовательно.

Ещё одной системой, поддающейся моделированию, является вещество, находящееся под поршнем внутри сосуда с проводящими тепло стенками. В такой системе регулируемыми параметрами оказываются число частиц, давление и температура. Это NPT-ансамбль. Давление в молекулярной динамике определяется следующим образом:

\begin{equation}
    P = \frac{1}{dV}[\sum_{i=1}{N}(\frac{p_i^2}{m_i} + \vec{r_i} \vec{F_i}) - dV]
\end{equation}

Для реализации такого случая требуется баростат. Одним из наиболее распространенных баростатов является баростат Шиноды, который состоит из цепочки Нозе-Гувера. Изменение давления происходит за счет объема вычислительной ячейки. Такой баростат эффективен для систем с резкими изменениями плотности как, например, в системах с фазовыми переходами. Он также сохраняет фазовый объем.

Существует также возможность моделирования открытой системы, однако в рамках данной диссертации в таких моделях не возникло необходимости, потому они описаны не будут.

\section{Фазовые переходы 1-го рода и влияющие на них факторы – массообмен, теплообмен и поверхностное натяжение}

Фазовые переходы, то есть скачкообразные изменения свойств вещества при плавном изменении внешних параметров (температуры, давления), представляют собой фундаментальное явление в физике конденсированного состояния. С термодинамической точки зрения, фазовые переходы традиционно классифицируются по работам Ландау на переходы первого и второго рода. Критерием классификации является поведение термодинамических потенциалов, таких как свободная энергия Гиббса $G$ .

При фазовом переходе первого рода скачком изменяются первые производные термодинамического потенциала, которые соответствуют таким экстенсивным параметрам, как энтропия  $S = -\left(\frac{\partial G}{\partial T}\right)_P$  и объём  $V = \left(\frac{\partial G}{\partial P}\right)_T$. Это означает, что переход сопровождается выделением или поглощением скрытой теплоты (теплота плавления, испарения) и скачкообразным изменением плотности. Типичными примерами служат плавление, кипение и сублимация.

В отличие от этого, при фазовом переходе второго рода первые производные термодинамического потенциала изменяются непрерывно (нет скачка плотности и выделения скрытой теплоты), однако разрыв испытывают вторые производные, такие как теплоёмкость  $C_P$, сжимаемость и термический коэффициент расширения. Примером является переход в сверхпроводящее состояние.

Ключевой особенностью фазовых переходов первого рода является их склонность к затруднённому протеканию в метастабильных условиях. Поскольку сосуществующие фазы при таком переходе имеют различную плотность и структуру, возникновение новой фазы внутри исходной не может произойти одновременно во всём объёме. Вместо этого система проходит через образование локальных неустойчивостей — мелких зародышей или кластеров новой фазы. Этот начальный стадийный процесс, известный как нуклеация (зародышеобразование), и является кинетическим «бутылочным горлышком» перехода.

Процесс перехода между фазами можно описать следующей последовательностью:

1. Нуклеация: В переохлажденной или пересыщенной исходной фазе вследствие флуктуаций спонтанно образуются наноразмерные зародыши новой фазы.

2. Рост: Те зародыши, размер которых превысил некий критический порог, становятся устойчивыми и начинают расти, потребляя вещество из исходной фазы.

3. Коалесценция: Растущие домены новой фазы сливаются друг с другом, формируя общую структуру.

4. Завершение перехода: Новая фаза полностью замещает исходную.

Наиболее устоявшимся и фундаментальным подходом к описанию стадии нуклеации является классическая теория нуклеации (КТН), разработанная в работах Фольмера, Вебера, Беккера и Дёринга. В её основе лежит концепция критического зародыша — кластера новой фазы, находящегося в неустойчивом равновесии с исходной фазой. Его существование является следствием конкуренции двух факторов: объёмной свободной энергии, выигрышной при образовании более стабильной фазы, и поверхностной свободной энергии, связанной с созданием границы раздела фаз и являющейся энергетическим барьером для процесса.

Таким образом, КТН предоставляет количественную модель для расчёта скорости образования зародышей, что является центральной задачей при описании кинетики фазовых переходов первого рода.

\subsection{Классическая теория нуклеации}



\section{Поверхностные явления и их описание}

\subsection{Поверхность Гиббса. Адсорбция}

В данной работе рассмотрены фазовые переходы для двух различных систем: слоя жидкости, окруженного паром, и уединенного кластера, на который налетают атомы. Поверхностные явления исследуются для первого случая с использованием теории Гиббса. Поверхностью в рамках данной теории является плоская поверхность с нулевой толщиной. Потому считается, что ячейка разделена на объёмы, соответствующие жидкой и газовой фазам:

\begin{equation}
    V = V_{liq} + V_{gas}
\end{equation}

Рассматриваются равновесные случаи, то есть выполняются все три условия равновесия: равенство давлений, температур и химических потенциалов. Для числа частиц при этом требуется учесть поверхностный избыток, который и показывает величину адсорбции:

\begin{equation}
    N = N_{liq} + N_{gas} + N_{exc}
\end{equation}

Здесь $N_{exc}$ \--- избыточное количество вещества. Поверхностная плотность избыточного количества вещества называется адсорбцией:

\begin{equation}
    \Gamma = N_{exc} / S,
\end{equation}

где $S$ \--- площадь поверхности.  

Выбор положения поверхности, очевидно, изменяет значение адсорбции для каждого из компонентов. Однако, в секции с результатами будет показано, что качественно картина не меняется. То есть изотерма адсорбции описывается уравнением того же вида независимо от способа задания поверхности и её положения. Одним из наиболее удобных способов задания поверхности является эквимолярная поверхность \--- поверхность, для которой $N_{exc} = 0$. Однако, в случае наличия нескольких компонентов это равенство не может выполняться для всех компонентов и для смеси одновременно. Следовательно, в случае смеси приходится выбирать либо компонент, для которого адсорбция будет равна нулю, либо выбирать поверхность так, чтобы суммарная адсорбция была нулевой.

Теоретические статьи по адсорбции? 

Однако, плоская поверхность \--- предполагаемый случай, а реальный профиль плотности выглядит иначе. Наиболее распространенные способы аппроксимации \--- гиперболический тангенс и функция ошибок:

\begin{align}
\rho(x) = \rho_g + \frac{\rho_l-\rho_g}{2}(tanh(\frac{x-x_0}{D})+1),
\end{align}

\begin{align}
\rho(x) = \rho_g + \frac{\rho_l-\rho_g}{2}(erf(\frac{x-x_0}{D})+1),
\end{align}

где $rho_g$ и $rho_l$ \--- плотности пара и жидкости соответсвенно, $x_0$ \--- положение эквимолярной поверхности, $D$ \--- харатерная толщина поверхности.

Визуально можно наблюдать адсорбцию на зависимости плотности от координаты, ось которой направлена перпендикулярно рассматриваемой поверхности. Она выглядит как пик по сравнению с плотностями в обеих фазах. 

Расчеты и эксперименты по адсорбции

\subsection{Поверхностное натяжение}

Одной из других интересных характеристик поверхности является поверхностное натяжение. Оно показывает велчичну свободной энергии, необходимую для увеличения площади поверхности на одну единицу:

\begin{equation}
    \sigma = (\frac{\partial F}{\partial S})_{N, V, T} = (\frac{\partial G}{\partial S})_{N, P, T},
\end{equation}

Статьи по поверхностному натяжению для относительно плоской поверхности.

\subsection{Нуклеация}

Поверхностное натяжение влияет на процесс нуклеации. Классическая теория нуклеации.

Статьи по нуклеации.


\section{Процессы переноса}